\documentclass[UTF8,a4paper,11pt]{ctexart}

\usepackage[a4paper,margin=2.05cm,top=2.0cm,bottom=2.1cm,headheight=24pt,headsep=8pt,footskip=26pt]{geometry}
\usepackage{fontspec}
\usepackage{xcolor}
\usepackage{graphicx}
\usepackage{booktabs}
\usepackage{tabularx}
\usepackage{array}
\usepackage{longtable}
\usepackage{enumitem}
\usepackage{tikz}
\usepackage{titlesec}
\usepackage{fancyhdr}
\usepackage{lastpage}
\usepackage{hyperref}
\usepackage{amssymb}
\usepackage{setspace}
\newcommand{\tracking}[2]{#2}

\usetikzlibrary{arrows.meta,positioning,fit,calc,shapes.geometric}

\definecolor{navy}{HTML}{132D68}
\definecolor{blue}{HTML}{2379C9}
\definecolor{cyan}{HTML}{16A7C7}
\definecolor{sky}{HTML}{EAF5FC}
\definecolor{ink}{HTML}{20304A}
\definecolor{muted}{HTML}{61738D}
\definecolor{line}{HTML}{D8E4F0}
\definecolor{gold}{HTML}{F2B84B}
\definecolor{green}{HTML}{2E9C70}
\definecolor{red}{HTML}{C85D5D}

\setmainfont{Noto Sans SC}
\setsansfont{Noto Sans SC}
\setCJKmainfont{Noto Sans SC}
\setCJKsansfont{Noto Sans SC}
\setlength{\parindent}{2em}
\setlength{\parskip}{0.35em}
\linespread{1.18}
\hypersetup{colorlinks=true,linkcolor=blue,urlcolor=blue}

\pagestyle{fancy}
\fancyhf{}
\fancyhead[L]{\small\color{muted}星轨 AI 五育成长智能体}
\fancyhead[R]{\small\color{muted}2026 iCAN AI 应用创新挑战赛}
\fancyfoot[C]{\small\color{muted}应用方案 \textbar\ 第 \thepage\ 页 / 共 \pageref{LastPage} 页}
\renewcommand{\headrulewidth}{0.4pt}
\renewcommand{\headrule}{\hbox to\headwidth{\color{line}\leaders\hrule height \headrulewidth\hfill}}

\titleformat{\section}{\Large\bfseries\color{navy}}{\thesection}{0.7em}{}
\titleformat{\subsection}{\large\bfseries\color{blue}}{\thesubsection}{0.6em}{}
\titleformat{\subsubsection}{\normalsize\bfseries\color{ink}}{\thesubsubsection}{0.5em}{}
\titlespacing*{\section}{0pt}{1.0em}{0.45em}
\titlespacing*{\subsection}{0pt}{0.8em}{0.25em}
\titlespacing*{\subsubsection}{0pt}{0.5em}{0.15em}

\setlist[itemize]{leftmargin=2.1em,itemsep=0.2em,topsep=0.25em}
\setlist[enumerate]{leftmargin=2.35em,itemsep=0.25em,topsep=0.25em}

\newcolumntype{Y}{>{\raggedright\arraybackslash}X}
\newcommand{\tagbox}[1]{\colorbox{sky}{\textcolor{blue}{\textbf{#1}}}}
\newcommand{\key}[1]{\textcolor{blue}{\textbf{#1}}}
\newcommand{\note}[1]{\par\noindent\colorbox{sky}{\parbox{0.94\linewidth}{\textcolor{ink}{#1}}}\par}

\begin{document}

\thispagestyle{empty}
\begin{tikzpicture}[remember picture,overlay]
  \fill[navy] (current page.north west) rectangle ([yshift=-8.2cm]current page.north east);
  \fill[cyan] ([xshift=-3cm,yshift=-5.4cm]current page.north east) circle (4.2cm);
  \fill[blue!70!navy] ([xshift=1.5cm,yshift=-1.5cm]current page.south west) circle (3.3cm);
  \draw[white,opacity=.18,line width=0.8pt] ([xshift=-8cm,yshift=-2.0cm]current page.north east) ellipse (5.5cm and 1.8cm);
  \draw[white,opacity=.18,line width=0.8pt] ([xshift=-8cm,yshift=-2.0cm]current page.north east) ellipse (5.5cm and 1.0cm);
\end{tikzpicture}
\vspace*{1.5cm}
\begin{flushleft}
{\color{white}\small\bfseries\tracking{1.5}{2026 iCAN AI 应用创新挑战赛}}\\[1.2cm]
{\color{white}\fontsize{31}{38}\selectfont\bfseries 星轨 AI 五育成长智能体}\par
\vspace{0.35cm}
{\color{cyan!20!white}\fontsize{15}{23}\selectfont\bfseries 面向高校学生综合素质评价的\par\hspace{1em}智能材料审核、成长诊断与干预平台}
\end{flushleft}
\vfill
\begin{tabular}{@{}ll@{}}
\textcolor{muted}{参赛赛道} & 软件赛道 / 教育科研场景 \\
\textcolor{muted}{作品形态} & 网页应用 + AI 材料审核智能体 \\
\textcolor{muted}{提交材料} & 应用方案 \\
\textcolor{muted}{项目团队} & 星轨智评项目组 \\
\textcolor{muted}{编制日期} & 2026 年 9 月 23 日
\end{tabular}
\vspace{1.2cm}
\newpage

\tableofcontents
\newpage

\section{项目概述}

\subsection{项目摘要}

高校学生的成长证明材料，分散在纸质证书、图片、PDF、微信群、邮箱和个人网盘中。辅导员和班级管理员需要人工判断材料是否有效、属于哪一类五育成果、对应多少分值，再通过 Excel 或多个系统完成统计。学生最终看到的往往只是总分和排名，难以理解自己的优势、短板和下一步行动。

\textbf{星轨 AI 五育成长智能体}面向这一具体场景，将 AI 的非结构化材料理解能力与高校综合素质评价的规则引擎、审批流程和数据权限体系结合起来，形成“材料理解 - 规则匹配 - 人工复核 - 五育计分 - 成长干预 - 复盘反馈”的闭环。

项目不是用大模型直接替代教师评分，而是让 AI 负责材料理解、信息抽取、分类建议、风险提示和成长计划生成；最终高利害评价结果仍由学校规则和授权人员确认，从而兼顾效率、可解释性与责任边界。

\begin{center}
\resizebox{\linewidth}{!}{%
\begin{tikzpicture}[node distance=6mm and 6mm,>=Latex]
  \tikzstyle{box}=[rounded corners=3pt,draw=line,fill=white,minimum width=2.75cm,minimum height=0.78cm,align=center,font=\small\bfseries,text=ink]
  \node[box,fill=sky] (a) {成长材料\newline 图片 / PDF / 说明};
  \node[box,fill=blue!8,right=of a] (b) {AI 材料理解\newline 抽取 / 分类 / 风险};
  \node[box,fill=cyan!10,right=of b] (c) {规则与人工复核\newline 建议分值 / 确认};
  \node[box,fill=green!10,right=of c] (d) {五育评价\newline 计分 / 排名 / 公示};
  \node[box,fill=gold!16,below=of b] (e) {成长诊断\newline 识别短板};
  \node[box,fill=gold!16,right=of e] (f) {行动计划\newline 执行 / 复盘};
  \draw[->,thick,blue] (a)--(b)--(c)--(d);
  \draw[->,thick,blue] (d.south) |- (e.east);
  \draw[->,thick,blue] (e)--(f);
  \draw[->,thick,blue] (f.north) |- (a.south);
\end{tikzpicture}
}%
\end{center}

\subsection{核心价值}
\begin{tabularx}{\linewidth}{>{\bfseries\color{navy}}p{2.7cm}Y Y}
\toprule
价值层 & 解决的问题 & 形成的结果 \\
\midrule
效率价值 & 减少人工录入、材料翻找和重复核算 & 将一次材料处理从“人工查找”转为“AI 初审 + 人工确认” \\
评价价值 & 让成长证据进入统一、可追踪的评价流程 & 材料、规则、审批和分数之间形成证据链 \\
育人价值 & 从只展示结果转向解释原因和指导行动 & 每名学生获得基于自身数据的四周成长计划 \\
治理价值 & 避免模型直接替代高利害教育决策 & AI 建议、规则计算、人工复核和操作日志相互制约 \\
\bottomrule
\end{tabularx}

\section{项目背景与痛点问题}

\subsection{政策与场景背景}

新时代高校人才培养强调德智体美劳全面发展，学生评价需要从单一成绩、结果排名逐步走向多维度、过程性和发展性评价。与此同时，教育数字化建设已经积累了教务、学工、团委和课程等数据，但“成长材料如何理解”“不同成果如何归入五育”“评价结果如何反哺成长”仍然是实际工作中的难点。

本项目选择高校综合素质评价作为 AI 应用切口，原因在于它同时具备：
\begin{itemize}
  \item 真实而高频的材料处理需求；
  \item 明确的规则、角色和审批责任；
  \item 可形成从输入材料到结果反馈的完整闭环；
  \item AI 能够发挥作用，但又不能脱离人工和规则独立决策；
  \item 具备从单校试点扩展到多校配置化部署的空间。
\end{itemize}

\subsection{当前痛点}

\begin{longtable}{>{\bfseries\color{navy}}p{2.5cm}p{5.1cm}p{6.1cm}}
\toprule
痛点类型 & 现状表现 & 直接影响 \\
\midrule
采集断层 & 证书、活动证明和成果说明分散在图片、PDF、纸质文件和聊天工具中，格式不统一。 & 学生重复提交，管理员需要反复确认材料来源和完整性。 \\
理解断层 & 材料中的赛事等级、日期、证书编号和成果类型需要人工阅读，且不同人员理解口径可能不同。 & 审核耗时，容易出现漏填、错分、重复申报。 \\
核算断层 & 五育归类、单次上限、学期累计上限和学院自定义规则需要人工对照。 & Excel 多次复制粘贴，结果难以复核，出错后追溯成本高。 \\
反馈断层 & 学生常看到总分、排名和是否通过，却看不到结构化的成长建议。 & 评价停留在“给分”，没有有效转化为育人行动。 \\
治理断层 & 不同角色需要看到不同范围的数据，但传统表格难以持续保证权限边界。 & 敏感教育数据存在越权查看、误发和责任不清风险。 \\
\bottomrule
\end{longtable}

\subsection{问题本质}

这些问题不是简单的“缺一个报表”，而是三个环节没有连起来：
\begin{enumerate}
  \item \textbf{材料没有被理解：} 非结构化证据无法直接参与评价；
  \item \textbf{规则没有被执行：} 规则散落在通知、表格和个人经验中；
  \item \textbf{结果没有被使用：} 评价结果没有进一步指导学生成长。
\end{enumerate}

因此，项目将 AI 放在最有价值的位置：把材料转换成可审计的评价证据，并将评价数据转换成可执行的成长行动。

\section{需求分析与目标用户}

\subsection{目标用户}

\begin{tabularx}{\linewidth}{>{\bfseries\color{navy}}p{2.5cm}Y Y}
\toprule
用户群体 & 主要任务 & 核心诉求 \\
\midrule
学生 & 上传成果材料、提交加分申报、查看评价结果、执行成长计划 & 操作简单、结果透明、知道为什么以及下一步做什么 \\
班级管理员 & 处理本班材料、核验申报、反馈补充信息 & 快速识别材料重点，减少重复录入，保持审核口径一致 \\
辅导员 & 终审、关注班级差异、推动学生成长 & 看见风险和短板，能够基于数据开展精准指导 \\
院系 / 学校管理员 & 配置规则、管理组织、查看统计与报表 & 形成全局视图，保障权限、审计和评价制度落地 \\
\bottomrule
\end{tabularx}

\subsection{需求分层}

\begin{enumerate}
  \item \textbf{必需需求：} 材料上传、材料信息识别、五育模块匹配、建议分值、人工复核、申报提交和规则计分。
  \item \textbf{效率需求：} 自动提取标题、日期、证书编号和获奖等级；识别缺失字段；提示可能重复申报；减少审批人员在多个页面间切换。
  \item \textbf{治理需求：} 按角色和组织范围控制数据访问；保留 AI 分析结果、人工修改和最终审批记录；支持模型不可用时的安全降级。
  \item \textbf{育人需求：} 根据五育结构、申报活跃度和待处理事项，生成带目标、周期和复盘指标的成长计划。
  \item \textbf{扩展需求：} 通过 OpenAI 兼容接口接入不同模型，保持模型层可替换；未来扩展 PDF OCR、批量审核和校级知识库。
\end{enumerate}

\subsection{用户故事}

\note{学生：我有一张竞赛证书，不确定应该归到哪个模块、能申请多少分，希望系统先帮我读懂材料，再让我确认后提交。}
\note{管理员：我需要快速知道这份材料是否缺日期、等级或编号，AI 可以先给出依据和风险，但最终仍由我确认。}
\note{辅导员：我不只想看到学生排名，还想知道哪些维度长期偏弱，以及下个月可以安排什么具体行动。}

\section{产品方案与功能设计}

\subsection{产品定位}

星轨定位为“高校综合素质评价之外的 AI 证据理解与成长干预层”。它不替代学校现有教务、学工或审批制度，而是补齐从材料到评价、从评价到成长的智能化环节。

\subsection{功能全景}

\begin{center}
\begin{tikzpicture}[node distance=8mm and 10mm,>=Latex]
  \tikzstyle{module}=[rounded corners=4pt,draw=line,fill=white,minimum width=4.0cm,minimum height=1.05cm,align=center,text=ink]
  \node[module,fill=sky] (m1) {AI 材料理解中心\\\small 上传、识别、抽取、风险提示};
  \node[module,fill=blue!7,right=of m1] (m2) {申报审批中心\\\small 学生提交、管理员复核、辅导员终审};
  \node[module,fill=cyan!8,below=of m1] (m3) {五育核算中心\\\small 模块配置、分值计算、排名回写};
  \node[module,fill=gold!16,right=of m3] (m4) {成长驾驶舱\\\small 诊断、目标、行动、复盘};
  \node[draw=navy,rounded corners=6pt,fit=(m1)(m2)(m3)(m4),inner sep=12pt,label={[text=navy,font=\bfseries]above:星轨 AI 五育成长智能体}] (box) {};
  \draw[->,very thick,blue] (m1)--(m2);
  \draw[->,very thick,blue] (m2)--(m3);
  \draw[->,very thick,blue] (m3)--(m4);
  \draw[->,very thick,blue] (m4.north) |- (m1.south);
\end{tikzpicture}
\end{center}

\subsection{主要功能说明}

\begin{tabularx}{\linewidth}{>{\bfseries\color{navy}}p{3.1cm}Y}
\toprule
功能模块 & 功能说明 \\
\midrule
AI 材料审核 & 上传图片或 PDF；调用模型或本地规则完成成果标题、日期、证书编号、获奖等级、候选模块、建议分值、置信度、判定依据和风险提示的生成。 \\
申报流程 & 学生确认 AI 建议后，将材料和结构化结果带入原有申报表；系统支持草稿、提交、审批、退回补充、终审和申诉等状态。 \\
规则核算 & 按启用模块、单次上限、学期累计上限和学院配置的五育映射计算结果，避免模型直接修改最终分数。 \\
AI 成长计划 & 根据学生五育分布、历史申报和待处理事项识别优先提升方向，生成四周目标、行动任务、证据要求和复盘节点。 \\
权限与审计 & 学生只能分析本人上传材料；管理员按班级、院系或学校范围访问；AI 建议、人工确认和操作日志可追踪。 \\
\bottomrule
\end{tabularx}

\newpage
\section{AI 核心作用与技术方案}

\subsection{AI 在系统中的核心作用}

本项目的 AI 不是一个孤立的聊天框，而是嵌入业务流程、能够调用工具并接受人工复核的应用智能体。它承担五项核心工作：
\begin{enumerate}
  \item \textbf{看懂材料：} 对图片材料进行视觉理解，结合文件名、标题和补充说明识别成果事实；
  \item \textbf{抽取字段：} 输出标题、日期、证书编号、等级和成果说明等结构化字段；
  \item \textbf{匹配规则：} 在启用的评价模块中选择候选模块，参考模块上限给出建议分值；
  \item \textbf{识别风险：} 标注缺日期、无法判断等级、材料不充分和需要人工核验等风险；
  \item \textbf{推动成长：} 从评价数据中识别短板，生成可执行、可复盘的行动计划。
\end{enumerate}

\subsection{智能体工作流}

\begin{center}
\begin{tikzpicture}[node distance=5mm and 8mm,>=Latex]
  \tikzstyle{flow}=[rounded corners=3pt,draw=blue!35,fill=sky,minimum width=3.2cm,minimum height=0.72cm,align=center,font=\small]
  \node[flow] (a) {用户上传材料};
  \node[flow,right=of a] (b) {文件安全校验};
  \node[flow,right=of b] (c) {模型 / 规则路由};
  \node[flow,below=of c] (d) {字段和事实抽取};
  \node[flow,left=of d] (e) {模块与分值建议};
  \node[flow,left=of e] (f) {风险与置信度};
  \node[flow,below=of f] (g) {人工确认};
  \node[flow,right=of g] (h) {申报与规则计分};
  \draw[->,thick,blue] (a)--(b)--(c)--(d)--(e)--(f)--(g)--(h);
\end{tikzpicture}
\end{center}

\subsection{模型层与降级机制}

模型层采用 OpenAI 兼容接口设计，可配置服务地址、API 密钥、模型名称和超时时间。当前系统支持两种运行模式：
\begin{itemize}
  \item \textbf{模型分析模式：} 图片材料以多模态请求发送给配置的模型，要求模型只返回约定 JSON 结构，后端再次进行字段归一化、分值上限裁剪和风险数组清洗。
  \item \textbf{本地规则模式：} 未配置密钥、模型超时或服务返回异常时，系统根据文件名、标题和说明中的关键词完成初步分类，明确标记“规则降级”和“需要人工复核”。
\end{itemize}

这一设计保证了演示和业务流程在模型不可用时仍可运行，同时避免把低置信度结果伪装成模型结论。当前图片材料支持视觉模型分析；PDF 已纳入上传和流程，但完整 PDF 文本 OCR 仍属于后续增强项。

\subsection{技术架构}

\begin{center}
\begin{tikzpicture}[node distance=6mm,>=Latex]
  \tikzstyle{layer}=[rounded corners=4pt,draw=line,minimum width=12.8cm,minimum height=0.82cm,align=center]
  \node[layer,fill=sky] (l1) {展示层：Vue 3 + Vite + Element Plus + ECharts};
  \node[layer,fill=blue!7,below=of l1] (l2) {业务层：Express REST API + 申报审批 + 五育核算 + 成长计划};
  \node[layer,fill=cyan!8,below=of l2] (l3) {智能层：AI 材料审核服务 + OpenAI 兼容模型 + 本地规则降级};
  \node[layer,fill=gold!16,below=of l3] (l4) {数据层：PostgreSQL + Sequelize + 上传文件 + 审计日志};
  \node[layer,fill=green!10,below=of l4] (l5) {治理层：六级角色权限、数据作用域、人工复核、操作留痕};
  \draw[->,very thick,navy] (l1)--(l2)--(l3)--(l4)--(l5);
\end{tikzpicture}
\end{center}

\subsection{开发工具与工程环境}

\begin{tabularx}{\linewidth}{>{\bfseries\color{navy}}p{3.0cm}Y}
\toprule
类别 & 使用工具 / 技术 \\
\midrule
前端开发 & Vue 3、Vite、Element Plus、Pinia、ECharts、Axios、Sass \\
后端开发 & Node.js、Express、Sequelize、JWT、Multer、Sharp \\
数据存储 & PostgreSQL；模型同时保留 MySQL 兼容设计和初始化脚本 \\
AI 接入 & OpenAI 兼容 Chat Completions 接口；支持可配置模型和服务地址 \\
文件处理 & 图片上传、压缩、静态文件安全访问；后续扩展 PDF OCR \\
工程管理 & Git、GitHub 分支备份、前端生产构建、Node 语法检查、降级模式冒烟测试 \\
\bottomrule
\end{tabularx}

\section{作品功能详述}

\subsection{AI 材料审核中心}

页面分为“上传与补充材料”和“AI 审核建议”两个工作区。左侧完成材料选择、模块偏好、标题和补充说明；右侧按审核结果展示置信度、建议模块、建议分值、获奖等级、识别日期、证书编号、判定依据和风险提示。

AI 结果不会直接写入最终成绩。用户必须点击“带结果去申报”，在申报表中再次确认模块、标题、日期、分值和附件后提交。这样既减少填写工作，也保留用户的知情和纠错机会。

\subsection{申报审批与五育核算}

原有业务流程作为可信底座：
\begin{enumerate}
  \item 学生选择评价模块，填写成果信息并上传证明材料；
  \item 系统校验申报学期、模块状态、重复申报和分值上限；
  \item 班级管理员进行材料初审，可通过、驳回或退回补充；
  \item 辅导员对符合条件的结果进行终审；
  \item 系统依据规则将分值归入五育模块，并回写汇总、排名和公示数据；
  \item 所有关键操作记录在审批记录和操作日志中。
\end{enumerate}

\subsection{AI 五育成长计划}

成长计划不是泛化鸡汤，而是根据当前学生数据生成：
\begin{itemize}
  \item 当前总分、五育分类分布和优先提升维度；
  \item 已通过申报、待处理申报和最近参与情况；
  \item 目标分值与可验证成长记录要求；
  \item 四周行动任务，例如参加志愿服务、完成学习复盘、参加体育活动或整理证明材料；
  \item 每周复盘节点，用于检查行动是否开始、证据是否完整和下一周期目标是否需要调整。
\end{itemize}

\subsection{权限与数据安全}

系统按照学生、班级管理员、辅导员、院系管理员、学校管理员和超级管理员进行分层。AI 材料接口额外校验文件地址和用户归属，普通用户只能分析本人上传的文件；组织管理员只能在授权的数据作用域中查看业务数据。

在教育评价场景中，安全不只意味着接口鉴权，还包括：模型不直接决定最终成绩、输出显示依据和置信度、低置信度结果必须人工确认、模型调用失败自动降级、生产密钥不进入代码仓库。

\newpage
\section{使用说明与交互指南}

\subsection{学生操作流程}

\begin{enumerate}
  \item 登录系统后进入“AI 材料审核”；
  \item 点击“选择图片或 PDF”，上传竞赛证书、志愿服务证明或活动材料；
  \item 可选择“优先匹配模块”，也可以交给 AI 自动判断；
  \item 补充成果标题和说明，点击“开始 AI 分析”；
  \item 阅读建议模块、建议分值、置信度、判定依据和风险提示；
  \item 对结果满意时点击“带结果去申报”；
  \item 在申报页确认学期、标题、日期、等级、分值和证明材料；
  \item 提交申报或保存草稿，在“我的申报”中追踪审批状态；
  \item 回到学生数据看板查看 AI 五育成长计划。
\end{enumerate}

\subsection{管理员操作流程}

\begin{enumerate}
  \item 登录后进入“考评审核中心”；
  \item 查看学生提交的材料、AI 建议和人工填写内容；
  \item 对照学校评分标准核验材料原件、日期、等级和成果真实性；
  \item 选择通过、退回补充或驳回，并填写必要意见；
  \item 需要时修改最终分值，系统执行模块上限和累计上限校验；
  \item 审批完成后查看五育计分、排名和操作记录。
\end{enumerate}

\subsection{交互原则}

\begin{tabularx}{\linewidth}{>{\bfseries\color{navy}}p{3.0cm}Y}
\toprule
交互原则 & 具体设计 \\
\midrule
先理解后提交 & AI 结果先展示，用户确认后才进入正式申报，不自动提交。 \\
结果可解释 & 每个建议配套置信度、判定依据和风险提示，避免只给一个结论。 \\
失败可恢复 & 模型调用失败时进入规则降级模式；用户仍可手动选择模块并继续流程。 \\
高利害人工把关 & AI 输出被定义为审核建议，最终评分由规则和授权管理员确认。 \\
行动有闭环 & 评价结果不止停留在看板，直接生成目标、任务和复盘节点。 \\
\bottomrule
\end{tabularx}

\subsection{典型演示场景}

\begin{enumerate}
  \item 上传“全国大学生创新创业大赛国家级获奖证书”图片；
  \item AI 识别“国家级”关键词，推荐创新创业或竞赛相关模块，给出建议分值；
  \item 结果页显示证书日期缺失或需要人工核验等风险；
  \item 用户将结果带入申报页，补充日期和说明后提交；
  \item 管理员确认后，系统自动计入五育评价；
  \item 学生看板识别其创新创业维度变化，并生成下一阶段成长计划。
\end{enumerate}

\section{创新性、实用价值与展示方案}

\subsection{创新点}

\begin{enumerate}
  \item \textbf{从材料理解切入五育评价：} 将 AI 视觉理解能力用于真实教育材料，而不是停留在通用问答。
  \item \textbf{AI 建议与规则治理结合：} 大模型负责非结构化理解，规则引擎负责分值边界，人工负责最终责任确认。
  \item \textbf{从评价结果走向成长干预：} 输出可执行的四周行动和复盘指标，使评价服务于育人。
  \item \textbf{权限感知与证据链：} 将数据作用域、人工复核和操作留痕内置到 AI 应用中。
  \item \textbf{可替换、可降级：} 通过 OpenAI 兼容接口接入不同模型，模型不可用时仍能安全运行。
\end{enumerate}

\subsection{与传统方案对比}

\begin{tabularx}{\linewidth}{>{\bfseries\color{navy}}p{3.0cm}Y Y Y}
\toprule
对比维度 & Excel + 群聊 & 普通管理系统 & 星轨 AI 五育成长智能体 \\
\midrule
材料处理 & 人工查看和转录 & 多数仅保存附件 & AI 抽取字段、分类和风险提示 \\
评分逻辑 & 分散在表格和经验中 & 规则可配置但缺少材料理解 & AI 建议 + 规则校验 + 人工确认 \\
结果反馈 & 总分和排名 & 总分、排名和报表 & 五育诊断、目标、行动和复盘 \\
安全治理 & 依赖人工管理 & 角色权限 & 作用域权限、归属校验和审计留痕 \\
模型故障 & 不适用 & 不适用 & 明确降级为本地规则并提示人工复核 \\
\bottomrule
\end{tabularx}

\subsection{与大赛评审标准的对应关系}

\begin{tabularx}{\linewidth}{>{\bfseries\color{navy}}p{3.0cm}Y}
\toprule
评审维度 & 项目对应内容 \\
\midrule
创新性 & 将多模态材料理解、规则治理和成长干预结合到高校综合素质评价闭环，形成与传统表格和普通管理系统的差异化。 \\
技术实现 & 采用 Vue 3、Express、OpenAI 兼容模型、本地规则降级、权限校验和审计留痕，支持在线演示和持续扩展。 \\
实用价值 & 面向真实的材料审核、申报审批、五育核算和学生成长指导需求，具有从院系试点扩展到校级部署的路径。 \\
用户体验 & 学生先获得 AI 解释再提交，管理员按证据和风险复核，模型失败时仍可继续人工处理。 \\
展示效果 & 方案提供完整操作流程和 5 分钟演示脚本，重点展示 AI 接收请求、分析决策、输出建议和协同审批。 \\
\bottomrule
\end{tabularx}

\subsection{5 分钟演示视频脚本}

\begin{tabularx}{\linewidth}{>{\bfseries\color{blue}}p{2.1cm}Y Y}
\toprule
时间 & 画面 & 讲解重点 \\
\midrule
0:00-0:30 & 痛点与产品定位 & 材料分散、人工审核、结果缺少成长建议。 \\
0:30-1:20 & 进入 AI 材料审核中心并上传图片 & 展示 AI 不是聊天框，而是嵌入申报流程的智能体。 \\
1:20-2:20 & 展示字段、模块、分值、置信度和风险 & 强调可解释输出和人工复核边界。 \\
2:20-3:10 & 带结果去申报并提交 & 展示 AI 模块与原有审批、规则引擎协同。 \\
3:10-4:20 & 管理员审核和五育计分 & 展示权限、上限校验、审批留痕和结果回写。 \\
4:20-5:00 & 学生看板成长计划 & 展示短板诊断、四周行动和复盘指标，完成价值闭环。 \\
\bottomrule
\end{tabularx}

\section{应用前景与商业模式}

\subsection{应用前景}

项目可以从高校学生综合素质评价切入，逐步扩展到：
\begin{itemize}
  \item 奖助学金和荣誉评审材料的智能预审；
  \item 第二课堂、志愿服务、创新创业和社会实践成果归档；
  \item 辅导员学生工作台和重点学生成长风险提示；
  \item 学院人才培养质量分析和专业实践活动评价；
  \item 高职、本科、研究生等不同培养层次的成长档案管理。
\end{itemize}

项目的扩展边界不是简单复制页面，而是将“材料类型、评分规则、组织权限和成长指标”做成可配置项，适应不同学校的制度差异。

\subsection{商业模式}

\begin{tabularx}{\linewidth}{>{\bfseries\color{navy}}p{3.0cm}Y}
\toprule
模式 & 说明 \\
\midrule
校级 SaaS 订阅 & 面向学校提供组织管理、规则配置、AI 材料审核、报表和运维服务，按学校规模或活跃用户数计费。 \\
院系试点包 & 面向单个学院或学生工作部门提供低门槛试点，先验证材料处理效率和学生使用率，再扩展到全校。 \\
私有化部署 & 对教育数据敏感、已有基础设施的学校提供本地化部署、模型适配、数据治理和技术支持。 \\
增值服务 & 提供规则迁移、历史材料整理、指标体系设计、AI 工作坊和使用培训。 \\
\bottomrule
\end{tabularx}

\subsection{落地路径}

\begin{enumerate}
  \item 选择一个学院或一个年级进行材料审核和五育申报试点；
  \item 用脱敏样本建立材料类型、模块分类和字段抽取的人工标注集；
  \item 比较 AI 辅助前后的处理耗时、人工采纳率和错误类型；
  \item 根据学校制度配置模块上限、审批流程和数据范围；
  \item 形成校级推广方案，接入更多成长记录和学生工作场景。
\end{enumerate}

\section{实施计划与评测方案}

\subsection{实施计划}

\begin{tabularx}{\linewidth}{>{\bfseries\color{blue}}p{2.8cm}Y Y}
\toprule
阶段 & 主要工作 & 交付物 \\
\midrule
当前版本 & AI 图片材料审核、申报结果带入、规则降级、成长计划和权限校验 & 可运行网页应用、演示账号和基础测试脚本 \\
短期增强 & PDF 文本 OCR、批量材料审核、人工纠错回流、模型调用统计 & OCR 服务、审核队列、评测报表 \\
中期扩展 & 学校规则知识库、更多材料类型、管理员配置中心 & 配置化规则和多场景模板 \\
长期发展 & 多校部署、模型微调或专用分类模型、成长效果评估 & 行业解决方案和标准化产品包 \\
\bottomrule
\end{tabularx}

\subsection{评测指标}

为避免只展示“模型能不能回答”，项目建立面向业务的评测指标：
\begin{itemize}
  \item \textbf{材料类型准确率：} AI 是否正确识别证书、志愿服务证明、竞赛材料等类型；
  \item \textbf{字段抽取准确率：} 标题、日期、证书编号和等级的字段准确程度；
  \item \textbf{五育模块分类准确率：} 与人工标注结果的一致程度；
  \item \textbf{人工采纳率：} 管理员直接采纳 AI 建议或仅做小幅修改的比例；
  \item \textbf{平均响应耗时：} 从上传完成到返回审核建议的时间；
  \item \textbf{降级率与可用性：} 模型不可用时，系统是否保持安全、可解释、可继续操作。
\end{itemize}

当前版本已提供降级模式冒烟测试；正式参赛材料应使用 20-50 份脱敏材料建立人工标注表，再填入实际测量结果，不使用未经验证的数字。

\subsection{风险与应对}

\begin{tabularx}{\linewidth}{>{\bfseries\color{navy}\raggedright\arraybackslash}p{3.4cm}Y Y}
\toprule
风险 & 可能影响 & 应对措施 \\
\midrule
模型幻觉或误识别 & 错误抽取、错误推荐模块 & 低温度 JSON 输出、字段归一化、置信度、风险提示和人工复核。 \\
材料隐私泄露 & 学生证书和身份信息暴露 & 权限校验、文件路径校验、生产密钥不入库、脱敏评测。 \\
评分责任不清 & AI 建议被误认为最终结果 & 明确 AI 只提供建议，最终计分由规则和授权人员完成。 \\
模型服务不可用 & 现场演示或业务流程中断 & 本地规则降级、明确标识降级状态、支持手动继续申报。 \\
规则变化 & 建议分值与学校制度不一致 & 模块上限和评价配置可维护，模型不直接写最终成绩。 \\
\bottomrule
\end{tabularx}

\section{结语}

星轨 AI 五育成长智能体的核心不是把一个大模型接入校园系统，而是围绕高校真实的综合素质评价场景，建立一条可运行、可解释、可复核、可持续改进的 AI 应用链路。

项目以材料理解解决“证据进不来”，以规则和审批解决“结果不可信”，以成长计划解决“评价用不上”。在教育评价这一高责任场景中，项目坚持“AI 提效但不越权，模型建议但不替代制度，自动分析但保留人工责任”的设计原则，具备软件赛道要求的在线演示能力、行业应用价值和持续扩展空间。

\end{document}
